\chapter{Credit Default Swaps}\label{ch:m2:credit-default-swaps}

Lehman Brothers filed for bankruptcy on 15 September 2008. \omterm{def:m2:credit-default-swaps:cds}{Credit default swaps}
on it with a gross notional of about USD~72 billion were registered, and only about a
third of that protection could have been bought to hedge Lehman bonds; the rest was
held by investors who had no bonds to hand over. On 10 October fourteen dealers
ran a two-stage auction on behalf of 358 parties; it priced Lehman's senior bonds at
8.625 per 100, and every \omterm{def:m2:credit-default-swaps:cds}{protection seller} that had signed up paid 91.375 per 100 of
notional, in cash. Offsetting positions shrank the money that changed hands to about
USD~5.2 billion. This chapter explains the contract that made those payments: its
standard form since 2009, how it is quoted and priced with a \omterm{def:m2:credit-default-swaps:hazard}{hazard rate}, who decides
that a default has happened, how the auction sets the price, and why the cheapest bond
decides the payout.

\section{The standard contract}

\begin{definition}[Credit default swap, protection buyer and seller, reference entity]\label{def:m2:credit-default-swaps:cds}
A \emph{credit default swap}\index{credit default swap} (CDS) is a bilateral contract in
which the \emph{protection buyer}\index{protection buyer} pays a regular premium on a
notional amount until maturity or until a \omterm{def:m2:credit-default-swaps:event}{credit event} of a named borrower, the
\emph{reference entity}\index{reference entity}, whichever comes first; after a \omterm{def:m2:credit-default-swaps:event}{credit
event} the \emph{protection seller}\index{protection seller} pays the buyer the notional
times one minus the recovery value of the entity's debt, fixed by settlement.
\end{definition}

A CDS is insurance on a borrower's debt, with two differences that matter: the buyer
need not own the debt, and the contract can be traded and offset like any swap. The
buyer is short the credit and the seller long, so a seller of protection has the same
exposure to default as a bondholder without having to fund the bond. On USD~10 million
of protection at 100 basis points a year, the buyer pays USD~25\,000 a quarter; if the
debt is worth 40 per 100 after default, the seller pays USD~6 million
(\cref{fig:m2:credit-default-swaps:flows}).

\begin{omfigure}
\centering
\begin{tikzpicture}[font=\small,
  box/.style={draw,thick,rounded corners=2pt,align=center,minimum height=1.1cm,minimum width=3.2cm,inner sep=3pt},
  arr/.style={-{Stealth[length=2.5mm]},thick}]
\node[box,draw=omDef,fill=omDef!8] (b) at (0,0) {protection buyer\\(short the credit)};
\node[box,draw=omThm,fill=omThm!8] (s) at (9.4,0) {protection seller\\(long the credit)};
\node[box,draw=gray,fill=gray!8] (r) at (4.7,-3.2) {reference entity\\(its bonds and loans)};
\draw[arr] ([yshift=3mm]b.east) -- node[above=2pt,font=\footnotesize,align=center] {premium: $c \times$ notional a year,\\quarterly, until default or maturity} ([yshift=3mm]s.west);
\draw[arr] ([yshift=-3mm]s.west) -- node[below=2pt,font=\footnotesize,align=center] {after a credit event:\\$(100 - \text{final price})\%$ of notional} ([yshift=-3mm]b.east);
\draw[dashed,thick,gray] (r.north) -- (4.7,-1.45);
\node[font=\footnotesize,text=gray,align=right,anchor=east] at ([xshift=-4pt]r.west) {credit event\\decided by the DC};
\node[font=\footnotesize,text=gray,align=left,anchor=west] at ([xshift=4pt]r.east) {auction sets\\the final price};
\end{tikzpicture}

\omcaption{The cash flows of a \omterm{def:m2:credit-default-swaps:cds}{credit default swap}. The buyer pays a running premium
until a \omterm{def:m2:credit-default-swaps:event}{credit event} or maturity; after a \omterm{def:m2:credit-default-swaps:event}{credit event}, determined by a committee, the
seller pays the notional times one minus the price of the entity's debt set by an
auction. Schematic.}\label{fig:m2:credit-default-swaps:flows}
\end{omfigure}

\begin{definition}[Standard coupon, upfront payment]\label{def:m2:credit-default-swaps:std}
A \emph{standard coupon}\index{standard coupon} is a fixed running premium, the same
for all contracts of a class, which the market trades instead of a coupon negotiated
for each trade. The \emph{upfront payment}\index{upfront payment} is the amount, in
per cent of notional, that one party pays the other when the trade starts so that the
contract with the standard coupon is fair at the market's spread; the \omterm{def:m2:credit-default-swaps:cds}{protection
buyer} pays it when the spread is above the coupon and receives it when below.
\end{definition}

Until 2009 each trade carried its own coupon, the par spread of the day, so two
offsetting trades done a week apart left two streams of different premiums that
could not be netted or cleared. From 8 April 2009, contracts on North American
companies moved to standard terms: a fixed coupon of 100 basis points for
investment-grade names and 500 for high yield, with the difference paid upfront;
maturities on the quarterly roll dates, 20 March, June, September and December;
premiums accruing from the last roll date before the trade, the stub settled with the
upfront; and restructuring no longer a \omterm{def:m2:credit-default-swaps:event}{credit event} for these contracts. The same change created
the committees and hard-wired the auctions of \cref{sec:m2:credit-default-swaps:events}.
Contracts on the same name and maturity became identical, and so could be netted
and, later, cleared by a central counterparty.

\begin{dated}{2026-09}{The size of the CDS market}\label{dat:m2:credit-default-swaps:size}
At the end of 2025, by the BIS statistics as summarised by ISDA: notional outstanding
USD~11.0 trillion, up 21.7\% in a year, of which USD~4.6 trillion single-name and
USD~6.4 trillion on several names (indices and tranches, next chapter); gross market
value USD~273.7 billion; 70.2\% of the notional, USD~7.7 trillion, cleared by central
counterparties.
\end{dated}

\section{Upfront and running}

To convert a spread into an \omterm{def:m2:credit-default-swaps:std}{upfront payment}, and to price a contract at all, one
needs a model of when the entity defaults.

\begin{definition}[Hazard rate, survival probability]\label{def:m2:credit-default-swaps:hazard}
The \emph{hazard rate}\index{hazard rate} $\lambda(t)$ of a \omterm{def:m2:credit-default-swaps:cds}{reference entity} is its
instantaneous rate of default at time $t$ given survival to $t$: the probability of
default in $[t, t + dt]$, given no default before $t$, is $\lambda(t)\,dt$. The
\emph{survival probability}\index{survival probability} to $T$ is $Q(T) =
\exp\bigl(-\int_0^T \lambda(u)\,du\bigr)$, which is $e^{-\lambda T}$ for a flat hazard
rate.
\end{definition}

A CDS has two legs. The premium leg pays $c$ a year, quarterly, while the entity
survives, plus the premium accrued up to a default; its value per unit of coupon is
the risky \emph{annuity} $A$ (chapter 9's annuity, with each payment weighted by the
probability of surviving to it). The protection leg pays $1 - R$ at default, with $R$
the \omterm{def:m2:corporate-bonds:covenant}{recovery rate} of chapter 21. With a flat rate $r$ and a flat hazard $\lambda$,
paying defaults at the end of the quarter in which they occur and counting half a
quarter of accrued premium,
\[
A = \sum_{k=1}^{4T} \tfrac14\, e^{-r t_k}\Bigl(Q(t_k) + \tfrac12\bigl(Q(t_{k-1}) - Q(t_k)\bigr)\Bigr),
\qquad
P = (1-R)\sum_{k=1}^{4T} e^{-r t_k}\bigl(Q(t_{k-1}) - Q(t_k)\bigr),
\]
with $t_k = k/4$. The par spread is $P/A$, the running premium that makes the contract
worth zero.

\begin{proposition}[The credit triangle]\label{prop:m2:credit-default-swaps:triangle}
With a flat \omterm{def:m2:credit-default-swaps:hazard}{hazard rate} $\lambda$, a constant recovery $R$ and premiums paid
continuously, the par spread is
\[
s = \lambda\,(1 - R),
\]
whatever the interest rate and the maturity.
\end{proposition}

\begin{proof}
With premiums paid continuously, the premium leg is worth $s\int_0^T e^{-(r+\lambda)t}\,dt$
and the protection leg $(1-R)\int_0^T \lambda\,e^{-(r+\lambda)t}\,dt$, since default
occurs in $[t, t+dt]$ with probability $\lambda e^{-\lambda t}dt$. The integrals
differ only by the constant factor $\lambda$, so the legs are equal when $s =
\lambda(1-R)$.
\end{proof}

With quarterly payments the triangle is no longer exact, but close: at a spread of
100 basis points and a recovery of 40\%, the \omterm{def:m2:credit-default-swaps:hazard}{hazard rate} that prices the five-year
contract is 1.6667\%, against $0.01/0.6 = 1.6667\%$ from the triangle, and the
five-year \omterm{def:m2:credit-default-swaps:hazard}{survival probability} is 92.0\%.

\begin{proposition}[Upfront from a quoted spread]\label{prop:m2:credit-default-swaps:upfront}
If a contract with \omterm{def:m2:credit-default-swaps:std}{standard coupon} $c$ is quoted at a spread $s$, and $\lambda_s$ is
the flat \omterm{def:m2:credit-default-swaps:hazard}{hazard rate} at which $s$ is the par spread, the \omterm{def:m2:credit-default-swaps:std}{upfront payment} per unit
notional, paid by the \omterm{def:m2:credit-default-swaps:cds}{protection buyer}, is
\[
U = (s - c)\,A(\lambda_s).
\]
\end{proposition}

\begin{proof}
At $\lambda_s$ the protection leg equals $s\,A(\lambda_s)$, by the definition of the
par spread. The contract with coupon $c$ is worth $P - cA = (s - c)A$ to the buyer,
which the buyer pays upfront to make the trade fair.
\end{proof}

\begin{example}[Two quotes, two conventions]\label{ex:m2:credit-default-swaps:upfront}
With a flat rate of 4\% and a recovery of 40\%, the five-year risky annuity is 4.332
at a spread of 100 basis points and 4.165 at 200. An investment-grade name quoted at
200 basis points, on the 100 basis point coupon, costs the buyer $1\% \times 4.165 =
4.16\%$ upfront, USD~416\,000 on USD~10 million. A high-yield name quoted at 300
basis points on the 500 coupon pays the buyer 8.01\% upfront; at 800 the buyer pays
9.98\%, and at 1\,200, 20.28\%. Investment-grade names are quoted in spread,
high-yield names in points upfront.
\end{example}

\Cref{fig:m2:credit-default-swaps:upfront} shows the upfront against the spread for
both coupons: it crosses zero at the coupon, and it bends because the annuity shrinks
as the \omterm{def:m2:credit-default-swaps:hazard}{hazard rate} rises. The slope near a spread is the annuity, so USD~10 million
of five-year protection near 100 basis points gains about USD~4\,300 for each basis
point the spread widens. \Cref{fig:m2:credit-default-swaps:survival} shows what the
spreads mean for survival.

\begin{omfigure}
\begin{tikzpicture}
\begin{axis}[width=11cm,height=5cm,
  xlabel={quoted spread (bp)},ylabel={upfront paid by the buyer (\%)},
  tick label style={font=\small},label style={font=\small},
  xmin=0,xmax=1500,ymin=-25,ymax=40,grid=major,grid style={gray!25},
  legend columns=2,legend style={font=\footnotesize,at={(0.5,-0.3)},anchor=north,draw=none,fill=none,/tikz/every even column/.append style={column sep=8pt}},
  table/col sep=comma]
\addplot[omDef,very thick] table[x=spread,y=c100]{figdata/markets-2/23-credit-default-swaps/upfront.csv};
\addplot[omThm,very thick] table[x=spread,y=c500]{figdata/markets-2/23-credit-default-swaps/upfront.csv};
\addplot[gray,thin,forget plot] coordinates {(0,0) (1500,0)};
\addplot[only marks,mark=*,mark size=1.6pt,black,forget plot] coordinates {(200,4.16) (300,-8.01) (800,9.98) (1200,20.28)};
\legend{100 bp coupon,500 bp coupon}
\end{axis}
\end{tikzpicture}

\omcaption{\omterm{def:m2:credit-default-swaps:std}{Upfront payment} of a five-year contract, in per cent of notional, against
the quoted spread, for the two \omterm{def:m2:credit-default-swaps:std}{standard coupons} (flat rate 4\%, recovery 40\%). The
upfront is zero where the spread equals the coupon, and negative, paid to the buyer,
below it; the dots are \cref{ex:m2:credit-default-swaps:upfront}. Illustrative; data:
the chapter's tutorial.}\label{fig:m2:credit-default-swaps:upfront}
\end{omfigure}

\begin{omfigure}
\begin{tikzpicture}
\begin{axis}[width=11cm,height=4.8cm,
  xlabel={years},ylabel={survival probability},
  tick label style={font=\small},label style={font=\small},
  xmin=0,xmax=10,ymin=0.2,ymax=1.02,grid=major,grid style={gray!25},
  yticklabel style={/pgf/number format/.cd,fixed,precision=1},
  legend columns=3,legend style={font=\footnotesize,at={(0.5,-0.3)},anchor=north,draw=none,fill=none,/tikz/every even column/.append style={column sep=8pt}},
  table/col sep=comma]
\addplot[omDef,very thick] table[x=t,y=s100]{figdata/markets-2/23-credit-default-swaps/survival.csv};
\addplot[omThm,very thick] table[x=t,y=s300]{figdata/markets-2/23-credit-default-swaps/survival.csv};
\addplot[omProp,very thick] table[x=t,y=s800]{figdata/markets-2/23-credit-default-swaps/survival.csv};
\legend{100 bp,300 bp,800 bp}
\end{axis}
\end{tikzpicture}

\omcaption{Survival probabilities implied by five-year spreads of 100, 300 and 800
basis points with a flat \omterm{def:m2:credit-default-swaps:hazard}{hazard rate} and a recovery of 40\%. The \omterm{def:m2:credit-default-swaps:hazard}{hazard rates} are
close to spread divided by $1 - R$: about 1.7\%, 5\% and 13\% a year. These are
risk-neutral probabilities, which include a premium for bearing default risk.
Illustrative; data: the chapter's tutorial.}\label{fig:m2:credit-default-swaps:survival}
\end{omfigure}

The probabilities are those that price the contract, not forecasts. Chapter 21
split a bond's spread into expected loss and a premium for risk and illiquidity;
the same split holds here, so a \omterm{def:m2:credit-default-swaps:hazard}{hazard rate} read from a spread is higher than the
default rate history would suggest.

\section{Credit events and determinations committees}\label{sec:m2:credit-default-swaps:events}

\begin{definition}[Credit event, determinations committee]\label{def:m2:credit-default-swaps:event}
A \emph{credit event}\index{credit event} is one of the events listed in a CDS
contract whose occurrence triggers the protection payment, such as the \omterm{def:m2:credit-default-swaps:cds}{reference
entity}'s bankruptcy or failure to pay. A \emph{determinations
committee}\index{determinations committee} (DC) is the body of dealers and investors
that decides, for all standard contracts at once, whether and when a credit event
has occurred and how it is settled.
\end{definition}

ISDA's disclosure material lists the \omterm{def:m2:credit-default-swaps:event}{credit events} a transaction may specify:
bankruptcy, failure to pay, restructuring, obligation acceleration, obligation default
and repudiation or moratorium; the contract says which apply. Bankruptcy and failure to
pay are the common ones for companies, restructuring matters for banks and sovereigns,
and whether a given event qualifies is often a question of interpretation. Before
2009 each pair of counterparties decided for itself, and two offsetting trades could
be triggered differently.

The 2009 reform created a committee for each of five regions (the Americas, Asia
excluding Japan, Japan, Australia and New Zealand, and Europe, the Middle East and
Africa), with, at its creation, ten dealers and five other members voting and ISDA as
secretary. Anyone party to a covered trade may ask a question; the committee decides
whether a \omterm{def:m2:credit-default-swaps:event}{credit event} has occurred and when, whether to hold an auction, which
obligations may be delivered, and whether the entity has a successor. Most decisions
need 80\% of the members, failing which an external panel decides, and they bind all
covered trades. A 60-day look-back replaced each trade's own start date, so that
offsetting trades done on different days are triggered by the same events.

The rules evolve with the cases. In 2013 the Dutch state nationalised the bank SNS
and expropriated its subordinated bonds, an event the definitions of the time did not
clearly cover. The 2014 definitions, in force from 22 September 2014, added a
\emph{governmental intervention} \omterm{def:m2:credit-default-swaps:event}{credit event} for financial entities outside the
United States and let \omterm{def:m2:credit-default-swaps:cds}{protection buyers} deliver whatever assets a bail-in or
restructuring turns the bonds into.

\section{Settlement auctions}

After a \omterm{def:m2:credit-default-swaps:event}{credit event}, the payment is $1 - R$ per unit of notional, and someone must
set $R$. Physical settlement, the original method, has the buyer deliver bonds of the
entity and receive par; but when most protection is held by investors without bonds, as
with Lehman, those buyers must buy bonds to deliver, which pushes their price up and the payout down.
Cash settlement needs one agreed price for bonds that may barely trade. The auction
provides it.

\begin{definition}[Auction final price]\label{def:m2:credit-default-swaps:final}
The \emph{auction final price}\index{auction final price} of a \omterm{def:m2:credit-default-swaps:event}{credit event} is the
price per 100 of par of the entity's deliverable debt, set by the settlement
auction, at which every CDS covered by the auction is cash settled and every
physical settlement request made in the auction is executed.
\end{definition}

The auction has two stages. In the first, each participating dealer submits a bid
and an offer, no more than 2 apart, for a fixed quotation size of bonds, and
physical settlement requests, for itself and its clients, to buy or sell bonds at the
final price, each no larger than the party's CDS position and in the same direction.
The administrators remove crossing pairs, highest bid against lowest offer, until no
bid reaches an offer, and average the better half of the remaining bids and offers,
rounded to the nearest eighth: the \emph{inside market midpoint}. The net of the
requests is the \emph{open interest}. A dealer whose quote crossed on the wrong side
of the midpoint, given the direction of the open interest, pays a penalty.

In the second stage, participants submit limit orders on the other side of the open
interest; the orders are filled best price first until the open interest is met, and
the last order filled sets the final price. The price may not be more than 1 above the
midpoint when the open interest is to sell, or 1 below when it is to buy.

\begin{omfigure}
\begin{tikzpicture}
\begin{axis}[width=11cm,height=5cm,
  xlabel={dealer (sorted by bid)},ylabel={price per 100},
  tick label style={font=\small},label style={font=\small},
  xmin=0.3,xmax=14.7,ymin=7.5,ymax=12.5,grid=major,grid style={gray!25},
  xtick={1,2,3,4,5,6,7,8,9,10,11,12,13,14},
  legend columns=4,legend style={font=\footnotesize,at={(0.5,-0.3)},anchor=north,draw=none,fill=none,/tikz/every even column/.append style={column sep=8pt}},
  table/col sep=comma]
\addplot[only marks,mark=triangle*,mark size=2.6pt,omDef] table[x=dealer,y=bid]{figdata/markets-2/23-credit-default-swaps/lehman.csv};
\addplot[only marks,mark=triangle*,mark options={rotate=180},mark size=2.6pt,omThm] table[x=dealer,y=offer]{figdata/markets-2/23-credit-default-swaps/lehman.csv};
\addplot[omProp,very thick,dashed] coordinates {(0.3,9.75) (14.7,9.75)};
\addplot[black,very thick] coordinates {(0.3,8.625) (14.7,8.625)};
\legend{bid,offer,midpoint 9.75,final 8.625}
\end{axis}
\end{tikzpicture}

\omcaption{The first stage of the Lehman Brothers auction, 10 October 2008: each of the
fourteen dealers' bid and offer for USD~5 million of bonds. One pair crossed (a bid and
an offer of 10); the average of the better half of the rest gave an inside market
midpoint of 9.75. The open interest was USD~4.92 billion to sell, and the second stage
set the final price at 8.625. Data: Helwege, Maurer, Sarkar and Wang, FRBNY Staff
Report 372, Box 1.}\label{fig:m2:credit-default-swaps:lehman}
\end{omfigure}

\begin{example}[Lehman Brothers]\label{ex:m2:credit-default-swaps:lehman}
In the first stage, the highest bid, 10, met an offer of 10 and both were removed; the
better seven of the remaining thirteen bids and offers averaged 9.804, rounded to
9.75. Requests to sell bonds totalled USD~5.69 billion and to buy USD~0.77 billion:
an open interest of USD~4.92 billion to sell, capping the final price at 10.75. The
dealer whose bid of 10 crossed, above the midpoint while the open interest was to
sell, paid $5\,\text{million} \times (10 - 9.75)\% = \text{USD}~12\,500$. In the second
stage, 453 bids to buy were submitted, from 10.75 down to 0.125; bids at 8.625 or
higher added up to the open interest, and the final price was 8.625
(\cref{fig:m2:credit-default-swaps:lehman}).
\end{example}

The auction does two things at once. It sets a single price, so an investor hedged
with a single-name contract against an index contract is not left with different
recoveries on each. And it turns physical settlement into a trade of the net: a
\omterm{def:m2:credit-default-swaps:cds}{protection buyer} that owns bonds and wants to be rid of them asks to sell them at the
final price and ends with par, while everyone else settles in cash. For Lehman, the
USD~72 billion of gross protection became USD~5.2 billion of net payments, because most
dealers had bought and sold protection on it.

The price is not guaranteed to match the bond market. Lehman's bonds traded at 34 on
the day of the bankruptcy, 13 the day before the auction and about 10 on the day;
the final price, pushed down by the large open interest to sell, was a little lower. In the Fannie Mae auction of 2008, strong demand in the
first stage pushed the final price of subordinated debt above that of the senior debt.

\section{The cheapest-to-deliver option}

A CDS protects not one bond but a class of the entity's debt: any obligation of the
required seniority, currency and form is deliverable. After a default, bonds of the
same seniority are claims on the same estate and should trade at similar prices, but
not always: a long bond with a low coupon and a short one with a high coupon can
differ before the bankruptcy settles, and in a restructuring the obligations may be
treated differently. A \omterm{def:m2:credit-default-swaps:cds}{protection buyer} that settles physically delivers the cheapest
of them, and in the auction the sellers of bonds deliver the cheapest too, so the
final price is set by the cheapest deliverable: the \omterm{def:m2:bond-futures:ctd}{cheapest-to-deliver} option of
chapter 6, now in credit. The option is worth most when deliverables differ in price,
which is when it matters most: restructurings that leave some bonds long-dated and
low-coupon, or bail-ins that convert bonds into other assets.

\section{Tutorial: pricing a CDS and replaying an auction}\label{tut:m2:credit-default-swaps:cds}

\begin{tutorial}
\textbf{Goal.} Price a CDS with a flat \omterm{def:m2:credit-default-swaps:hazard}{hazard rate}, convert quoted spreads into \omterm{def:m2:credit-default-swaps:std}{upfront
payments} for both \omterm{def:m2:credit-default-swaps:std}{standard coupons}, check the credit triangle, and replay the Lehman
auction.
\textbf{End state:} \cref{fig:m2:credit-default-swaps:upfront,fig:m2:credit-default-swaps:survival,fig:m2:credit-default-swaps:lehman,fig:m2:credit-default-swaps:orders},
\cref{ex:m2:credit-default-swaps:upfront,ex:m2:credit-default-swaps:lehman} and the
numbers of the weekend problem.
\begin{tutsteps}
\item \textbf{Pricing}: the two legs, par spread, upfront, and the \omterm{def:m2:credit-default-swaps:hazard}{hazard rate} implied by
a spread.
\omcode{code/firm/cds/firm_cds.py}{14}{53}{A CDS with a flat hazard rate: legs, par spread, upfront and the implied hazard.}\label{lst:m2:credit-default-swaps:price}
\item \textbf{The auction}: inside market midpoint, open interest, final price and cash
settlement.
\omcode{code/firm/cds/firm_cds.py}{58}{96}{The two-stage settlement auction.}\label{lst:m2:credit-default-swaps:auction}
\item \textbf{Run} \texttt{cds\_demo.tutorial()}, \texttt{cds\_demo.problem()} and
\texttt{fig\_cds.py}; build and test the C++20 and Rust versions in
\texttt{code/firm/cds/cpp} and \texttt{code/firm/cds/rust}.
\end{tutsteps}
\textbf{What to change next.} Replace the flat hazard by a curve bootstrapped from one-,
three-, five- and seven-year spreads, and see how much the five-year upfront changes;
then let the \omterm{def:m2:corporate-bonds:covenant}{recovery rate} vary and check that the upfront barely moves when the
\omterm{def:m2:credit-default-swaps:hazard}{hazard rate} is recalibrated to the same spread.
\end{tutorial}

\section{Build: the CDS pricer}\label{bld:m2:credit-default-swaps:cds}

\begin{build}
\bfield{Purpose} The miniature firm quotes and marks CDS on its bond issuers: upfront
and running conversions, \omterm{def:m2:credit-default-swaps:hazard}{hazard rates} for its credit models, and the auction's
arithmetic when a name defaults. Like the bond calculator of chapter 3, it exists in
Python, C++20 and Rust, with the same tests.
\bfield{Interface} \texttt{Cds(maturity, coupon, recovery, freq)};
\texttt{legs(cds, lam, r)}; \texttt{par\_spread}; \texttt{upfront};
\texttt{hazard\_from\_spread(cds, spread, r)}; \texttt{upfront\_from\_spread};
\texttt{inside\_market\_midpoint(bids, offers)}; \texttt{open\_interest(requests)};
\texttt{final\_price(imm, oi, orders)}; \texttt{cash\_settlement(notional, price)}.
\bfield{Rules} Flat hazard and flat continuously compounded rate; quarterly premiums
with half a period of accrued premium on default; defaults paid at the end of the
quarter; auction prices per 100 of par, midpoint rounded to an eighth, second-stage
price capped one point beyond the midpoint; the dealers' first-stage quotes are not
carried into the second stage.
\bfield{Acceptance tests} \texttt{code/firm/cds/tests/}, \texttt{cpp/firm\_cds\_test.cpp},
\texttt{rust/src/lib.rs}: the par spread round-trips through the implied hazard; the
upfront is zero at the coupon and changes sign across it; the Lehman first stage gives
9.75 and an open interest of 4\,920; the cap binds at 10.75.
\bfield{Stretch} Piecewise-constant hazard curves bootstrapped from a term structure
of spreads; exact dates and accrual on the roll-date calendar; a CDS on an index
(next chapter); the risk of a CDS book by name and tenor (One Quant Book 6).
\end{build}

\begin{omsources}
\item J.~Helwege, S.~Maurer, A.~Sarkar and Y.~Wang, ``Credit default swap auctions'',
Federal Reserve Bank of New York Staff Report 372, May 2009.
\item Paul, Weiss, ``The Big Bang Protocol and a new structural framework for credit
default swaps'', client memorandum, 24 March 2009.
\item F.~Carruzzo, ``New ISDA 2014 Credit Derivatives Definitions'', Harvard Law School
Forum on Corporate Governance, August 2014.
\item ISDA, Disclosure Annex for Credit Derivative Transactions.
\item ISDA, ``Key trends in the size and composition of OTC derivatives markets in the
second half of 2025'', July 2026.
\end{omsources}

\section{Exercises}

\begin{exercise}[$\star$]\label{exo:m2:credit-default-swaps:1}
An investor buys USD~10 million of protection at 100 basis points a year, paid
quarterly. What does it pay each quarter, and what does it receive if the entity
defaults and the \omterm{def:m2:credit-default-swaps:final}{auction final price} is 40?
\end{exercise}

\begin{exercise}[$\star$]\label{exo:m2:credit-default-swaps:2}
A name's five-year spread is 300 basis points and the recovery is taken as 40\%. Use
the credit triangle to estimate its \omterm{def:m2:credit-default-swaps:hazard}{hazard rate} and its five-year \omterm{def:m2:credit-default-swaps:hazard}{survival probability}.
\end{exercise}

\begin{exercise}[$\star$]\label{exo:m2:credit-default-swaps:3}
List four questions a \omterm{def:m2:credit-default-swaps:event}{determinations committee} decides after a possible \omterm{def:m2:credit-default-swaps:event}{credit event}.
\end{exercise}

\begin{exercise}[$\star\star$]\label{exo:m2:credit-default-swaps:4}
A five-year contract on an investment-grade name is quoted at 50 basis points. With the
chapter's assumptions, what is the upfront on USD~10 million, and who pays it?
\end{exercise}

\begin{exercise}[$\star\star$]\label{exo:m2:credit-default-swaps:5}
Why did \omterm{def:m2:credit-default-swaps:std}{standard coupons} make it possible to net and clear CDS, when contracts with
each trade's own par spread did not?
\end{exercise}

\begin{exercise}[$\star\star$]\label{exo:m2:credit-default-swaps:6}
After a default two bonds of the same seniority are deliverable, one trading at 30 and
one at 35. What does a \omterm{def:m2:credit-default-swaps:cds}{protection buyer} with USD~10 million of protection receive by
physical settlement, and which price should the auction reflect?
\end{exercise}

\begin{exercise}[$\star\star\star$]\label{exo:m2:credit-default-swaps:7}
\emph{Coding.} With \texttt{firm\_cds}, find the \omterm{def:m2:credit-default-swaps:hazard}{hazard rate} implied by five-year
spreads of 100, 300 and 800 basis points, and compare each with the credit triangle.
Explain the gap.
\end{exercise}

\begin{exercise}[$\star\star\star$]\label{exo:m2:credit-default-swaps:8}
\emph{Find the flaw.} ``The CDS spread is the market's probability of default per
year: a name at 300 basis points has a 3\% chance of defaulting each year.'' Correct it.
\end{exercise}

\section{Problem: The Auction}

\begin{problem}[{Weekend problem --- the two stages of a settlement auction}]\label{pb:m2:credit-default-swaps:1}
Replay the Lehman Brothers auction of 10 October 2008 with the first-stage submissions
of \cref{fig:m2:credit-default-swaps:lehman}: fourteen dealers' bids and offers for
USD~5 million of bonds, and physical settlement requests (USD millions) to buy 130, 612
and 30 and to sell 755, 870, 141, 480, 464, 1\,470, 390, 574, 191, 170 and 187. Only
the range and median of the second-stage orders were published; use the illustrative
book of limit bids to buy in \cref{fig:m2:credit-default-swaps:orders}, consistent with
the published result. Four illustrative participants settle: A, a pension fund, holds USD~20
million of the bonds and USD~20 million of protection; B, a hedge fund, bought USD~50
million of protection and holds no bonds; C, a dealer, sold USD~100 million of protection
and bought USD~60 million; D, an insurer, sold USD~30 million of protection.

\begin{omfigure}
\begin{tikzpicture}
\begin{axis}[width=11cm,height=4.8cm,
  xlabel={limit bid to buy (price per 100)},ylabel={cumulative size (USD bn)},
  tick label style={font=\small},label style={font=\small},
  xmin=7.3,xmax=11,ymin=0,ymax=14,grid=major,grid style={gray!25},
  x dir=reverse,
  legend columns=2,legend style={font=\footnotesize,at={(0.5,-0.3)},anchor=north,draw=none,fill=none,/tikz/every even column/.append style={column sep=8pt}},
  table/col sep=comma]
\addplot[omDef,very thick,const plot,mark=*,mark size=1.4pt] table[x=price,y expr=\thisrow{cum}/1000]{figdata/markets-2/23-credit-default-swaps/orders.csv};
\addplot[omThm,very thick,dashed] coordinates {(11,4.92) (7.3,4.92)};
\addplot[black,thick,forget plot] coordinates {(8.625,0) (8.625,14)};
\node[font=\footnotesize,anchor=east] at (axis cs:8.66,11.5) {final 8.625};
\legend{bids filled from the highest,open interest 4.92}
\end{axis}
\end{tikzpicture}

\omcaption{The second stage of the auction on an illustrative book of limit bids to buy
bonds: the bids are filled from the highest price down until they meet the open
interest of USD~4.92 billion to sell, which happens at 8.625, the published final
price. Bids below 7.5 are off the chart. Illustrative book; open interest and result:
FRBNY Staff Report 372.}\label{fig:m2:credit-default-swaps:orders}
\end{omfigure}

\medskip\noindent\textbf{Part I --- The first stage.}
\begin{enumerate}
\item Which bid and offer cross, and are removed?
\item Average the better half of the remaining bids and offers, and round it to the
inside market midpoint.
\item Compute the open interest and its direction.
\item Which dealer pays a penalty, and how much?
\item Why was the open interest to sell bonds?
\end{enumerate}

\medskip\noindent\textbf{Part II --- The second stage.}
\begin{enumerate}[resume]
\item Which orders are submitted, and what is the highest price allowed?
\item Fill the book from the highest bid and find the final price.
\item What would the final price be if all bids were at 12?
\item If only USD~1 billion had been bid at 8.625, what would the final price be?
\item By how much did the final price differ from the midpoint, and why?
\end{enumerate}

\medskip\noindent\textbf{Part III --- Settlement.}
\begin{enumerate}[resume]
\item A cash settles: what does it receive, and what is its bond worth at the final
price?
\item A instead asks to sell its USD~20 million of bonds in the auction: what does it
receive in all?
\item What does B receive?
\item What does C pay?
\item D asks to buy USD~30 million of bonds in the auction: what does it pay, and what
does it own afterwards?
\end{enumerate}

\medskip\noindent\textbf{Part IV --- Judgement.}
\begin{enumerate}[resume]
\item Why did USD~72 billion of gross protection turn into only USD~5.2 billion of
payments?
\item The bonds traded at about 10 on the day of the auction. Why could the final
price be lower?
\item What would have happened with physical settlement alone, with more protection
than bonds?
\item State the \emph{named result}: the final price and each participant's
settlement.
\item In one sentence: what does the auction add to a CDS?
\end{enumerate}
\end{problem}

\section{Interview questions}

\begin{interviewq}[$\star$ \iqroles{trader, researcher}]\label{iq:m2:credit-default-swaps:1}
What is a \omterm{def:m2:credit-default-swaps:cds}{credit default swap}? Who pays what, and when?
\end{interviewq}

\begin{interviewq}[$\star$ \iqroles{trader, bank}]\label{iq:m2:credit-default-swaps:2}
Why are CDS traded with a fixed coupon and an \omterm{def:m2:credit-default-swaps:std}{upfront payment}?
\end{interviewq}

\begin{interviewq}[$\star\star$ \iqroles{researcher}]\label{iq:m2:credit-default-swaps:3}
Derive the relation between a CDS spread, the \omterm{def:m2:credit-default-swaps:hazard}{hazard rate} and the \omterm{def:m2:corporate-bonds:covenant}{recovery rate}.
\end{interviewq}

\begin{interviewq}[$\star\star$ \iqroles{trader, researcher}]\label{iq:m2:credit-default-swaps:4}
How does a CDS settlement auction work, and why is it needed?
\end{interviewq}

\begin{interviewq}[$\star\star$ \iqroles{trader, researcher}]\label{iq:m2:credit-default-swaps:5}
What is the CDS--bond basis, and how would you trade a negative basis?
\end{interviewq}

\begin{interviewq}[$\star\star\star$ \iqroles{developer}]\label{iq:m2:credit-default-swaps:6}
Design a system that values a book of 50\,000 CDS positions every night and settles it
after a \omterm{def:m2:credit-default-swaps:event}{credit event}.
\end{interviewq}
